\clearpage
\appendix
\renewcommand{\thesubsection}{\thesection.\arabic{subsection}}
\section{支撑材料结构}

支撑材料按下表组织；数值归档、图形和官方结果工作簿分别保留，便于独立重建与核验。

\begin{longtable}{p{0.38\textwidth}p{0.52\textwidth}}
  \toprule
  文件路径 & 内容与用途 \\
  \midrule
  \endfirsthead
  \toprule
  文件路径 & 内容与用途 \\
  \midrule
  \endhead
  \texttt{result1.xlsx} & 问题一完整购电与充放电策略 \\
  \texttt{result2.xlsx} & 问题二全年计划及紧急购电结果 \\
  \texttt{result3.xlsx} & 问题三全年计划、调整及紧急购电结果 \\
  \texttt{result4-2.xlsx} & 波动电价下问题二的完整结果 \\
  \texttt{result4-3.xlsx} & 波动电价下问题三的完整结果 \\
  \texttt{code/} & 数据处理、模型求解、结果导出与独立校验程序 \\
  \texttt{figures/} & 正文图形及其可复现数据 \\
  \texttt{runs/problem1/} & 问题一各次运行配置、完整策略、敏感性及校验归档 \\
  \texttt{problem1-dashboard.html} & 问题一可离线查看的交互调度图 \\
  \texttt{runs/problem2/} & 问题二一月校准、逐日策略、对照、敏感性与校验归档 \\
  \texttt{problem2-dashboard.html} & 问题二日期选择与逐时离线交互图 \\
  \texttt{visualization-gallery.html} & 全文科学图形索引及PDF、SVG、PNG下载 \\
  \texttt{runs/visuals/} & 独立视觉版本、样式规范、源码及数值保护哈希 \\
  \texttt{issued\_plans.csv} & 问题二各起报日48小时预测及辅助规划，位于对应数值归档 \\
  \texttt{boundary/} & 24/48/72小时及远端残值的完整对照轨迹，位于问题二归档 \\
  \texttt{runs/problem3/} & 问题三一月规则检验、14种策略逐日轨迹、发布计划与调整记录、校验归档 \\
  \texttt{issued\_adjustments.csv} & 问题三每次调整后的剩余时段计划购电量、风险带与参考储能轨迹，位于对应归档 \\
  \texttt{runs/problem4/} & 问题四一月价格检验、8种策略逐日轨迹、价格预测误差与校验归档 \\
  \texttt{AI工具使用详情.pdf} & AI 工具使用过程、人工修改与核验记录 \\
  \bottomrule
\end{longtable}

\clearpage
\section{关键算法流程}

问题一执行流程为：校核附件与时间映射，构建并分层求解整数模型，计算线性松弛与对偶校验，独立审计运行轨迹，计算对照与敏感性，回读校验工作簿，最后将成功结果发布到固定入口并保留版本归档。运行失败时保存诊断，不发布为有效结果。可在项目根目录执行 \texttt{make solve-q1} 重算，执行 \texttt{make check-q1} 对最近归档作独立检查。

问题二执行流程为：以同一信息截止点生成多提前期预测，仅按已成熟的预测误差校准风险轮廓；一月选型后冻结参数，每日求解48小时代理规划，仅将前144时段确定为当天计划购电量，逐时执行并继承真实日末SOC；完成完整年度边界对照、局部敏感性及工作簿回读。运行\texttt{make solve-q2}重算，运行\texttt{make check-q2}核验实际及辅助轨迹，并因果重建保存的预测。

问题三执行流程为：读取附件3并按对齐检验将整点预报锚定到发布后对应起始时段，以发布前刚结束的实测时段作首锚点插值到10分钟；光伏点预报与本地趋势预测等权融合，负荷沿用问题二预测；每日0:00以同一48小时风险规划确定当天计划购电量，在6:00、12:00、18:00按最新预报、实际SOC与0:00计划重解剩余时段的有偿调整LP，只更新当日尚未执行时段的计划量；逐时贪心执行并继承真实SOC；按最终计划购电量与0:00计划量的差额一次结算。运行\texttt{make solve-q3}重算14种策略与一月检验，运行\texttt{make check-q3}独立核验轨迹、调整记录、因果重建与工作簿。

问题四执行流程为：读取附件4并采用与附件2相同的连续日历时标；以已完成日期的价格和本地净负荷预测生成因果日前价格预测（两周同日均值与净负荷回归等权融合），在调整时刻按已观测的当日价格作加性修正；分别以问题二信息集（计划固定）和问题三信息集（三时点有偿调整）回放全年，规划用预测价、账单按相应时段实时价结算；另以附件1固定轮廓、当日价格已知、无日内修正、仅计划为对照。运行\texttt{make solve-q4}重算，运行\texttt{make check-q4}核验结算价格、轨迹、调整记录、价格因果重建与两份工作簿。

图形采用统一语义颜色、子图编号与最终印刷字号，配置见\path{code/config/visual_style.json}，详细约束及设计依据见\path{code/vis/STYLE_GUIDE.md}。图形组织参考Nature科学图形指南与NeurIPS历史格式说明\cite{nature-figure-guide,neurips-figure-style}，论文版式仍以本竞赛规范为准。执行\texttt{make figures}可从已核验归档重绘，不重复求解；流程核对数值文件与工作簿哈希不变，并单独保存视觉版本。矢量图以160 mm全宽插入，区间阶梯保留10分钟分辨率；配对重采样区间与确定性敏感性连线明确区分。

\clearpage
\section{问题一参数与时间解释的补充分析}
\label{app:q1-sensitivity}

对容量、充放电功率及效率进行单因素重新优化，结果如图~\ref{fig:q1-sensitivity}。容量改变时同步保持 SOC 下限为额定容量的 $10\%$、上限为 $90\%$、初末电量为 $50\%$；因此，该组结果代表整套储能规格变化，不能解释为固定初始电量下的纯容量边际收益。

\begin{figure}[htbp]
  \centering
  \includegraphics[width=\textwidth]{q1_sensitivity.pdf}
  \caption{设备参数对最优费用的影响。a--c分别改变容量、功率上限与单程效率；各点均重新优化，实心点及竖虚线对应主设定。纵轴为局部费用范围，连线不表示统计置信区间。}
  \label{fig:q1-sensitivity}
\end{figure}

额定容量由 $6000$ 增至 $18000\ \mathrm{kWh}$ 时，最优费用由 $39314.00$ 降至 $30511.32$ 元；功率上限由 $2500$ 增至 $5000\ \mathrm{kW}$ 时，费用由 $35208.64$ 降至 $33801.48$ 元，再增至 $7500\ \mathrm{kW}$ 时为 $33715.70$ 元。这表明，在本日输入下，提高功率上限的边际节费趋于减弱，但该观察不能替代包含设备投资的容量选择分析。

主设定按国家标准把题面的 $90\%$ 读作往返效率并取 $\eta_c=\eta_d=\sqrt{0.9}$；若改读为双侧各 $0.9$（往返 $0.81$），费用升至 $35126.85$ 元，说明效率口径须在后续各问中保持一致。光伏输入分别缩放为原值的 $0.8$ 和 $1.2$ 倍并重新优化时，费用分别为 $40562.89$ 和 $29924.85$ 元；这些属于确定性输入敏感性，不能视为既定计划应对真实预测误差的表现。

将附件时间改作时段终点，并在同样初末电量条件下重新优化，费用为 $33801.50$ 元，与主解释相差约 $0.02$ 元。虽然全天费用接近，指定时段的映射并不相同，因此仍保存两套完整策略及逐项对应关系。主结果以模板时段起点解释为准，为后续各问提供统一的时间、设备及费用核算接口。

\clearpage
\section{问题二补充图形}
\label{app:q2-figures}

\begin{figure}[htbp]
 \centering\includegraphics[width=\textwidth]{../output/figures/q2_selected_days.pdf}
 \caption{指定日期的因果运行。a、c、e、g为0时计划购电量、实际净负荷与紧急补购，共用功率尺度；b、d、f、h为对应储电量，虚线为设备边界。实测曲线用于事后展示和当前反馈，不作为0时制定计划时的已知信息。}
 \label{fig:q2-days}
\end{figure}

\begin{figure}[htbp]
 \centering\includegraphics[width=\textwidth]{../output/figures/q2_boundary.pdf}
 \caption{跨日规划的运行机制。a：334天逐时平均储电量；b：两个时段的日均充电量，条形从零起。三方案预测器、分位数和控制器相同；平均曲线不代表任一实际单日轨迹。}
 \label{fig:q2-boundary}
\end{figure}

\clearpage
\section{问题三补充图形}
\label{app:q3-figures}

\begin{figure}[htbp]
 \centering\includegraphics[width=\textwidth]{../output/figures/q3_selected_days.pdf}
 \caption{问题三指定日期的因果运行。a、c、e、g：$0{:}00$计划购电量（点线）、最终计划购电量（实线）、实际净负荷与紧急购电，共用功率尺度，竖虚线为$6{:}00$、$12{:}00$、$18{:}00$调整时刻；b、d、f、h：对应储电量，虚线为设备边界。}
 \label{fig:q3-days}
\end{figure}

表~\ref{tab:q3-horizon}按发布时刻与提前期分段给出调整时刻风险轮廓的实际覆盖率与光伏点预报误差，并与冻结旧版信息链对照：同构残差使当日剩余段的覆盖率贴近名义值；次日覆盖段的光伏误差高于本地两日预测，但$0.8$分位数轮廓的覆盖率仍接近名义值。

\input{../output/tables/q3_horizon_layout.tex}

\clearpage
\section{问题四指定日期结果}
\label{app:q4-dates}

两族主策略在四个指定日期的结果按表1至表3格式给出；全天购电费均为按交付时段实时价结算的实际账单。

\begingroup
\renewcommand{\arraystretch}{0.9}
\input{../output/tables/q42_purchase.tex}
\input{../output/tables/q42_storage.tex}
\endgroup
\clearpage
\input{../output/tables/q42_emergency.tex}
\begingroup
\renewcommand{\arraystretch}{0.9}
\input{../output/tables/q43_purchase.tex}
\input{../output/tables/q43_storage.tex}
\endgroup
\clearpage
\input{../output/tables/q43_emergency.tex}

\clearpage
\section{补充推导与独立核验细节}
\label{app:validation-details}

\subsection{风险指标与经济分位数的边界}

为刻画极端补购风险，记日紧急购电费为 $\mathcal L$。可在目标中加入风险项 $\lambda\operatorname{CVaR}_\alpha(\mathcal L)$，其定义为
\begin{equation}
  \operatorname{CVaR}_\alpha(\mathcal L)
  =\min_{\zeta}\left\{\zeta+
  \frac{1}{1-\alpha}\mathbb E[(\mathcal L-\zeta)^+]\right\},
  \quad 0<\alpha<1.
  \label{eq:cvar}
\end{equation}
该项刻画高损失尾部，$\lambda$ 表示额外风险偏好。$5$ 倍紧急购电价格已经反映不足的经济代价，故期望费用模型本身即可作为主模型，CVaR 为可选扩展。实际费用、终端价值与风险项应分别报告，不能将优化目标值整体当作购电现金支出。


若允许需求准确揭示后以$1.5p$增购、以$0.5p$退款，单时段费用为
\[
 C(q)=pq+1.5p\mathbb E[(D-q)^+]-0.5p\mathbb E[(q-D)^+].
\]
最优0时计划购电量为中位数，说明后续调整机会会改变初始计划。随机价格的反例为：以90\%概率取$D=0,P=1$，以10\%概率取$D=100,P=10$。普通0.8分位数为零，期望费用500元；计划100 kWh的期望费用仅190元。故未知且与需求相关的价格不能从期望费用中约去。

\subsection{因果信息与辅助段核验}

\paragraph{问题2核验}
独立核验覆盖实际轨迹的守恒、SOC连续、设备边界、互斥、补购用途和费用；另检查每份288时段规划的参考物理可行性、预测截止日期、成熟残差边界及前144时段计划购电量一致性。测试修改未来实测，确认此前多提前期预测和执行前缀不变，并检验年末预测视窗完整。另以\texttt{make audit-causality}对已发布归档逐日复核：将当日及之后的全部附件2实测替换为随机数后重新预测与规划，334天的0时计划购电量均与归档一致（最大差$3\times10^{-12}$ kWh以内），规划目标恒等于视窗内计划购电费；同一复核对问题三、四的0时计划与规划价格同样通过。另以\texttt{make compare-time}固定已选参数比较两种时标：连续日历解释的334天现金费用比日循环解释高424.18元，紧急购电多516.18 kWh，指定日期的紧急购电结论不变。官方工作簿回读与48096条实际轨迹一致，第二天虚拟购电不进入提交结果。

\paragraph{问题3核验}
独立程序对\QThreeStrategyCount 条全年轨迹重算能量平衡、SOC转移与连续性、设备边界、充放电互斥、紧急电用途、计划购电取用恒等式、调整恒等式$q_t=q_t^0+\delta_t^+-\delta_t^-$及$\delta_t^+\delta_t^-=0$、计划版本不早于其生效时段、首个调整时刻前无调整，以及计划费、调整结算、紧急费与总费用的逐项恒等式，最大残差不超过$2\times10^{-5}$ kWh。全年\QThreeCash 元的现金费用由\QThreePlanCost、\QThreeAdjustCost、\QThreeEmergencyCost 万元三项精确加总。这些轨迹共求解\QThreeLPSolves 个线性规划（含一月影子检验全程为\QThreeLPSolvesTotal 个），最大原始--对偶目标差$\QThreeDualGap$元，被状态、残差或对偶证书拒绝的求解为\QThreeSolverFailures 次（任一拒绝都会中止回放，故完整归档的该计数必为零）。信息链四组变体的$0{:}00$计划288个时段与每份调整记录（主策略\QThreeAdjustmentRecords 次）的完整辅助段均已保存：辅助段满足状态转移、电量与功率边界，标记为未承诺，不进入任何账单；其点预报、风险轮廓与偏差修正由原始附件在归档配置下因果重建，最大差在$10^{-5}$ kWh以内，并对抽样记录在改写发布后观测、后续日期与未发布版本后重建第二次，结果逐项相同。测试用例改写尚未发布的版本与发布后的实测，确认此前的计划购电量与执行不变；两个信息链开关均关闭即为旧版信息链在当前口径下的同参对照；在旧版归档配置（双侧$0.9$效率、日循环时标）下它曾逐时复现2026年9月11日的冻结归档（最大差$10^{-12}$量级），口径变更后该复现断言按配置差异自动跳过并记录。权重取零且不调整时逐时复现问题二已发布策略。官方工作簿\texttt{result3.xlsx}的两张购电表分别为$0{:}00$计划及其费用、最终计划购电量及全天实际账单，充放电量与紧急购电量按最终执行填写，回读与轨迹一致。


\paragraph{问题4核验}
独立程序对\QFourStrategyCount 条全年轨迹重算物理、计划购电量、调整与费用恒等式，并逐时核对结算价格等于附件四相应时段的价格；两族$0{:}00$计划的288个时段与4-3每次调整的完整辅助段均已保存、通过物理与承诺标记核验，其规划价格、点预报与风险轮廓由附件四历史、附件三版本与本地预测在归档配置下因果重建，改写当日及以后的价格不改变任何一天的规划价格；附件一固定轮廓且无日内修正的价格模型逐时复现问题三的计划购电量与执行，两份工作簿回读与轨迹一致。



\FloatBarrier
\input{../output/tables/q3_information_chain.tex}
\FloatBarrier
\clearpage
\section{运行环境与复现入口}
\label{app:reproduction}

完整实现保存在支撑材料的\path{code/}目录。本附录保留补充证据、复现入口和四问关键实现；数据读写、绘图、导出及测试程序通过文件索引查阅。代码节选保留原文件行号，导入和依赖见对应源文件。

\subsection{环境、输入与配置}

Python依赖列于\path{code/requirements.txt}，主要使用NumPy与Pandas处理数据、SciPy求解规划、openpyxl读写工作簿、Matplotlib绘图。论文由XeLaTeX与latexmk编译。保留项目目录结构，在根目录运行以下命令。

原始附件位于\texttt{doc/Problem/}下的C题附件目录，按只读方式加载。设备与交易参数集中在\path{code/config/problem.json}；官方结果写入\path{output/results/}，不覆盖输入模板。

\subsection{命令与执行顺序}

\noindent\begin{tabularx}{\textwidth}{p{0.41\textwidth}L}
\toprule
根目录命令 & 用途 \\
\midrule
\texttt{make inspect} & 读取附件并检查输入结构。 \\
\texttt{make check-q1} 至 \texttt{make check-q4} & 依次检查四问最近成功归档，不重算优化。 \\
\texttt{make solve-q1} 至 \texttt{make solve-q4} & 需要重算时，按问题一至四顺序执行。 \\
\texttt{make test} & 运行手算案例、因果边界及接口测试。 \\
\texttt{make figures} & 从数值归档重绘程序生成的图形。 \\
\texttt{make check} & 编译论文并检查PDF结构。 \\
\texttt{make ai-details} & 单独构建AI工具使用详情。 \\
\bottomrule
\end{tabularx}

\subsection{复现时的核对重点}

先核验已保存结果，再按需要重算。四问独立检查分别覆盖物理约束、费用账本、预测信息边界与工作簿回读；辅助段只用于检验规划，其中的次日虚拟购电不形成正式计划，也不进入账单。

各问运行配置随数值版本保存，比较实验应使用同一版本内的预测、轨迹和账本。最终核对逐时记录、论文表格与五份工作簿一致；图一的独立图像资产另行维护，不由数值绘图命令重新生成。

\clearpage
\section{核心文件与结果索引}
\label{app:source-index}

下表路径相对于\path{code/}。模型模块位于\path{src/cumcm_c/}；各问工作流负责按顺序调用模型、保存诊断与发布结果。完整源码及测试随项目保留，论文不逐文件打印。

\begin{longtable}{p{0.57\textwidth}p{0.35\textwidth}}
\toprule
文件或目录 & 复核重点 \\
\midrule
\endfirsthead
\toprule
文件或目录 & 复核重点 \\
\midrule
\endhead
\path{config/problem.json} & 设备、交易与实验参数 \\
\path{src/cumcm_c/data.py} & 附件读取、时标与输入检查 \\
\path{src/cumcm_c/models/common.py} & 共享电池参数与轨迹接口 \\
\path{src/cumcm_c/models/problem1.py} & 整数模型、分层目标与证书 \\
\path{src/cumcm_c/models/forecasting.py} & 多提前期预测与成熟残差 \\
\path{src/cumcm_c/models/problem2.py} & 辅助规划、当天计划购电量与控制 \\
\path{src/cumcm_c/models/problem3.py} & 发布预报、风险带与调整 \\
\path{src/cumcm_c/models/problem4.py} & 因果价格与交付时段结算 \\
\path{src/cumcm_c/models/evaluation.py} & 问题一物理审计与基线 \\
\path{src/cumcm_c/models/problem2_evaluation.py} & 问题二实际轨迹独立核验 \\
\path{src/cumcm_c/models/problem3_evaluation.py} & 计划购电调整与费用独立核验 \\
\path{scripts/solve_problem1.py} 至 \path{scripts/solve_problem4.py} & 四问重算入口 \\
\path{scripts/validate_problem1.py} 至 \path{scripts/validate_problem4.py} & 归档独立核验入口 \\
\path{tests/}；\path{vis/} & 手算与因果测试；绘图源码 \\
\bottomrule
\end{longtable}

\subsection{从论文结论追溯到运行数据}

数值版本位于\path{output/runs/problem1/}至\path{output/runs/problem4/}。逐时记录区分预测计划、最终生效的计划购电量、实际执行及辅助段；验证时应在同一版本内比较，避免拼接不同配置的结果。

\noindent\begin{tabularx}{\textwidth}{p{0.29\textwidth}L}
\toprule
核验对象 & 应查看的材料与边界 \\
\midrule
计划购电与运行账单 & 实际轨迹、调整记录及五份结果工作簿；计划、调整和紧急费用分别核算。 \\
跨日规划与信息链 & \path{issued_plans.csv}、\path{issued_adjustments.csv}及对应归档配置；次日辅助段不形成正式计划、不结算。 \\
对照与结论 & 同参版本的逐日统计、敏感性与重采样区间；区间含零或风险升高的结果仍予保留。 \\
\bottomrule
\end{tabularx}

\clearpage
\section{关键实现节选}
\label{app:code-excerpts}

以下代码直接从现有源文件截取，行号对应原文件。每段展示一项核心机制，需在原模块及其依赖环境中使用，不作为独立脚本执行。

\begingroup
\lstset{basicstyle=\ttfamily\linespread{1}\fontsize{9}{11}\selectfont,numberstyle=\fontsize{7}{8}\selectfont\color{gray}}

\subsection{问题一：构造能量守恒与互斥约束}

\path{models/problem1.py}中的\texttt{build\_model}将购电、充放电、弃光和储电量组织为统一变量。两组等式分别对应能量平衡与状态转移；整数变量限制同一时段的充放电模式，边界固定日初、日末电量。

\lstinputlisting[language=Python,firstline=22,lastline=59,firstnumber=22,caption={确定性模型的变量、边界与约束}]{../code/src/cumcm_c/models/problem1.py}

分层求解及原始--对偶核验由同文件的\texttt{solve}继续完成。本节只展示可行域构造；求解状态、最优费用与独立审计结果见正文和对应数值归档。

\clearpage
\subsection{问题二：成熟残差与当天计划边界}

预测误差必须同时匹配起报日与提前期。下段按已完成目标日截取残差，再保留负荷、光伏成对误差向量，以此形成风险轮廓；该构造不声称给出了完整的跨日联合分布。

\lstinputlisting[language=Python,firstline=71,lastline=101,firstnumber=71,caption={按提前期构造已成熟误差与风险轮廓}]{../code/src/cumcm_c/models/forecasting.py}

48小时规划返回整段参考购电量与储电量，但仅将当天144时段确定为正式计划。次日参考状态也不替代实际执行后的储电量；下一天从真实末态继续运行。

\lstinputlisting[language=Python,firstline=99,lastline=102,firstnumber=99,caption={仅确定当天计划购电量}]{../code/src/cumcm_c/models/problem2.py}

\clearpage
\subsection{问题三：发布预报的跨日覆盖与残差重建}

每次发布的24小时预报可能覆盖次日的一部分。第一个函数仅在该版本覆盖范围内融合外部预报，范围之外沿用当天0:00已经生成的本地预测；第二个函数按同一规则重建历史预测，并排除目标日尚未结束的误差。

\lstinputlisting[language=Python,firstline=67,lastline=99,firstnumber=67,caption={已发布版本的次日辅助段与同构成熟残差}]{../code/src/cumcm_c/models/problem3.py}

这些预测进入辅助规划，不形成次日正式计划。当日有偿调整的线性规划、实际执行与最终结算分别见同文件的\texttt{adjust\_contract}和\texttt{run\_day}；调整量始终相对于0:00计划记账。

\clearpage
\subsection{问题四：价格预测与实际结算分离}

下列方法属于\path{models/problem4.py}中的\texttt{PriceModel}类。日前价格组合只使用已完成日期；在日内时刻$t$，加性修正只使用$t$之前已观测价格。\texttt{settle}另行返回交付日实价，防止将规划价格当作实际账单价格。

\lstinputlisting[language=Python,firstline=29,lastline=54,firstnumber=29,caption={因果价格预测的分量与等权组合}]{../code/src/cumcm_c/models/problem4.py}

\lstinputlisting[language=Python,firstline=72,lastline=80,firstnumber=72,caption={日内修正与交付时段结算接口}]{../code/src/cumcm_c/models/problem4.py}

类的构造函数及\texttt{today}、\texttt{tomorrow}接口见同文件。价格预知模式仅用于信息条件放宽的启发式对照，不作为最优费用下界；主策略仍采用因果预测模式。

\endgroup
